% ch5_e 第5章 §5.5–§5.6 (书页 239–249 / p254–p264)
%JOIN%
这使我们能够计算
\begin{align*}
&\sum_{i}\left(\left\langle c^{\dagger}_{\alpha i}c_{\alpha i}\right\rangle \left\langle c^{\dagger}_{\beta i}c_{\beta i}\right\rangle + \left\langle c^{\dagger}_{\alpha i}c_{\beta i}\right\rangle \left\langle c_{\alpha i}c^{\dagger}_{\beta i}\right\rangle\right)\\
&=\sum_{i}\left(n_{\uparrow i}+n_{\downarrow i}\right)^{2}+n_{\uparrow i}\left(1-n_{\uparrow i}\right)+n_{\downarrow i}\left(1-n_{\downarrow i}\right)\\
&=\sum_{i}n_{\uparrow i}+n_{\downarrow i}+2n_{\uparrow i}n_{\downarrow i}\\
&=N+2N^{-2}_{\mathrm{site}}\sum_{i}\sum^{\prime}_{\pmb{k},\pmb{k}'}\left(n^{-}_{\uparrow\pmb{k}}+n^{+}_{\uparrow\pmb{k}}+2(-)^{i}u_{\pmb{k}}v_{\pmb{k}}\left(n^{-}_{\uparrow\pmb{k}}-n^{+}_{\uparrow\pmb{k}}\right)\right)\times\\
&\qquad\qquad \left(n^{-}_{\downarrow\pmb{k}'}+n^{+}_{\downarrow\pmb{k}'}-2(-)^{i}u_{\pmb{k}'}v_{\pmb{k}'}\left(n^{-}_{\downarrow\pmb{k}'}-n^{+}_{\downarrow\pmb{k}'}\right)\right)\\
&=N+2N^{-1}_{\mathrm{site}}\sum^{\prime}_{\pmb{k},\pmb{k}'}\left(\left(n^{-}_{\uparrow\pmb{k}}+n^{+}_{\uparrow\pmb{k}}\right)\left(n^{-}_{\downarrow\pmb{k}'}+n^{+}_{\downarrow\pmb{k}'}\right)-\frac{B^{2}}{E_{\pmb{k}}E_{\pmb{k}'}}\left(n^{-}_{\uparrow\pmb{k}}-n^{+}_{\uparrow\pmb{k}}\right)\left(n^{-}_{\downarrow\pmb{k}'}-n^{+}_{\downarrow\pmb{k}'}\right)\right)
\end{align*}
利用这一结果，我们发现
\begin{align*}
\left\langle \Psi^{\mathrm{exc}}_{B}\right|H\left|\Psi^{\mathrm{exc}}_{B}\right\rangle&=\sum^{\prime}_{\pmb{k},\alpha}\left(u^{2}_{\pmb{k}}-v^{2}_{\pmb{k}}\right)\left(n^{-}_{\alpha\pmb{k}}-n^{+}_{\alpha\pmb{k}}\right)\xi_{\pmb{k}} \tag{5.5.6}\\
&\quad+\frac{U}{2}\sum_{i}\left(\left\langle c^{\dagger}_{\alpha i}c_{\alpha i}\right\rangle \left\langle c^{\dagger}_{\beta i}c_{\beta i}\right\rangle + \left\langle c^{\dagger}_{\alpha i}c_{\beta i}\right\rangle \left\langle c_{\alpha i}c^{\dagger}_{\beta i}\right\rangle\right)\\
&=-\sum^{\prime}_{\pmb{k},\alpha}\frac{\xi^{2}_{\pmb{k}}}{E_{\pmb{k}}}\left(n^{-}_{\alpha\pmb{k}}-n^{+}_{\alpha\pmb{k}}\right)+\frac{U}{N_{\mathrm{site}}}\sum^{\prime}_{\pmb{k},\pmb{k}'}\left(n^{-}_{\uparrow\pmb{k}}+n^{+}_{\uparrow\pmb{k}}\right)\left(n^{-}_{\downarrow\pmb{k}'}+n^{+}_{\downarrow\pmb{k}'}\right)\\
&\quad+\frac{U}{N_{\mathrm{site}}}\sum^{\prime}_{\pmb{k},\pmb{k}'}\frac{B^{2}}{E_{\pmb{k}}E_{\pmb{k}'}}\left(n^{-}_{\uparrow\pmb{k}}-n^{+}_{\uparrow\pmb{k}}\right)\left(n^{-}_{\downarrow\pmb{k}'}-n^{+}_{\downarrow\pmb{k}'}\right)+\frac{U}{2}N
\end{align*}
其中 $N=\sum_{i}n_{\uparrow i}+n_{\downarrow i}$ 是电子总数，而 $n_{\alpha i}$ 是格点 $i$ 上自旋为 $\alpha$ 的电子数。

把 $\langle\Psi^{\mathrm{exc}}_{B}|H|\Psi^{\mathrm{exc}}_{B}\rangle$ 展开到 $\delta n^{\pm}_{\alpha\pmb{k}}$ 的线性阶，我们得到
\begin{align*}
\left\langle \Psi^{\mathrm{exc}}_{B}\right|H\left|\Psi^{\mathrm{exc}}_{B}\right\rangle&=E_{\mathrm{ground}}+\frac{U}{2}\delta N-\sum^{\prime}_{\pmb{k},\alpha}\frac{\xi^{2}_{\pmb{k}}}{E_{\pmb{k}}}\left(\delta n^{-}_{\alpha\pmb{k}}-\delta n^{+}_{\alpha\pmb{k}}\right)\\
&\quad+\frac{U}{N_{\mathrm{site}}}\sum^{\prime}_{\pmb{k},\pmb{k}',\alpha}\left(\left(\delta n^{-}_{\alpha\pmb{k}}+\delta n^{+}_{\alpha\pmb{k}}\right)-\frac{B^{2}}{E_{\pmb{k}}E_{\pmb{k}'}}\left(\delta n^{-}_{\alpha\pmb{k}}-\delta n^{+}_{\alpha\pmb{k}}\right)\right)\\
&=E_{\mathrm{ground}}+\sum^{\prime}_{\pmb{k},\alpha}E_{\pmb{k}}\left(\delta n^{+}_{\alpha\pmb{k}}-\delta n^{-}_{\alpha\pmb{k}}\right)+U\delta N+O(\delta n^{2})
\end{align*}
其中我们用到了 $N=\sum^{\prime}_{\pmb{k},\alpha}(n^{-}_{\alpha\pmb{k}}+n^{+}_{\alpha\pmb{k}})$、$\sum^{\prime}_{\pmb{k}}1=N_{\mathrm{site}}/2$ 以及 $1=\frac{U}{N_{\mathrm{site}}}\sum^{\prime}_{\pmb{k}'}E^{-1}_{\pmb{k}'}$。

\section{非线性 $\sigma$ 模型}

\subsection{自旋密度波态的非线性 $\sigma$ 模型}

\begin{keypoints}
\item 非线性 $\sigma$ 模型描述对称性破缺基态附近集体涨落的动力学。
\item Hubbard--Stratonovich 变换。
\end{keypoints}

上一节所用的平均场方法未能再现 SDW 态中的无能隙自旋波激发。为了得到自旋波激发，我们必须超越平均场理论，比如采用 RPA。不过，这里我们将使用半经典方法以及与之相关的非线性 $\sigma$ 模型来研究自旋波激发。在半经典方法中，我们需要首先导出自旋波的低能有效拉格朗日量。

在我们的变分方法中，为了生成具有均匀反铁磁自旋极化 $M$ 的尝试波函数，必须在"变分"哈密顿量中引入一个 $B$ 场。把自旋波包括进来的一种自然方式，是在"变分"哈密顿量中用依赖于空间的 $\pmb{B}_{\pmb{i}}$ 场取代 $B$，这会在新的尝试波函数中生成依赖于空间的反铁磁自旋极化 $\pmb{M}_{\pmb{i}}$。新尝试波函数的能量是 $\pmb{M}_{\pmb{i}}$ 的泛函，即 $E[\pmb{M}_{\pmb{i}}]$，可以把它看作自旋波的能量。然而，$E[\pmb{M}_{\pmb{i}}]$ 是不依赖时间的自旋波涨落的能量，只是自旋波的势能。为了得到有效拉格朗日量，我们还需要依赖时间的自旋波涨落的动能。获得自旋波动能（以及势能）的最容易的方法，是使用 5.4.1 节中发展的路径积分方法。

为了得到自旋波的低能有效理论，我们注意到 Hubbard 项 $\frac{U}{2}c^{\dagger}_{\alpha i}c_{\alpha i}c^{\dagger}_{\beta i}c_{\beta i}$ 在 $n_{i}=0$ 时等于 0，在 $n_{i}=1$ 时等于 $U/2$，在 $n_{i}=2$ 时等于 $2U$；同时还注意到，$\pmb{S}^{2}_{i}$ 在 $n_{i}=0$ 时等于 0，在 $n_{i}=1$ 时等于 $3/4$，在 $n_{i}=2$ 时等于 0，其中
\begin{equation*}
\pmb{S}_{i}=\frac{1}{2}c^{\dagger}_{i}\sigma c_{i}
\end{equation*}
是格点 $i$ 上的总自旋算符。于是
\begin{equation*}
\frac{U}{2}c^{\dagger}_{\alpha i}c_{\alpha i}c^{\dagger}_{\beta i}c_{\beta i}=Un_{i}-\frac{2U}{3}\pmb{S}^{2}_{i}
\end{equation*}
我们将把 $Un_{i}$ 项吸收到化学势中，从而把它舍去。这样，Hubbard 模型就由下面的路径积分描述：
\begin{equation*}
Z=\int \mathcal{D}c^{*}\mathcal{D}c\;\mathrm{e}^{\mathrm{i}\int \mathrm{d}t\,(L_{0}+L_{I})}
\end{equation*}
其中
\begin{align*}
L_{0}&=\mathrm{i}\sum_{i}c^{*}_{\alpha i}\partial_{t}c_{\alpha i}-\sum_{\langle ij\rangle}-t(c^{*}_{\alpha i}c_{\alpha j}+c^{*}_{\alpha j}c_{\alpha i})\\
L_{I}&=\frac{U}{6}\sum_{i}c^{\dagger}_{i}\sigma c_{i}c^{\dagger}_{i}\sigma c_{i}
\end{align*}

这里的技巧是把相互作用项写成
\begin{equation*}
\mathrm{e}^{\mathrm{i}\int \mathrm{d}t\,\frac{U}{6}c^{\dagger}_{i}\sigma c_{i}c^{\dagger}_{i}\sigma c_{i}}=\text{constant}\times\int \mathcal{D}[B_{i}(t)]\;\mathrm{e}^{\mathrm{i}\int \mathrm{d}t\,(-)^{i}B_{i}(t)c^{\dagger}_{i}\sigma c_{i}-\frac{3}{2U}B^{2}_{i}(t)}
\end{equation*}
它把费米子的四次项变成了二次项。这样的变换称为 Hubbard--Stratonovich 变换。于是配分函数可以改写如下：
\begin{align*}
Z&=\int \mathcal{D}c^{*}\mathcal{D}c\,\mathcal{D}\pmb{B}\;\mathrm{e}^{\mathrm{i}\int \mathrm{d}t\,L_{B}-\mathrm{i}\int \mathrm{d}t\,\frac{3}{2U}\sum_{i}B^{2}_{i}} \tag{5.6.1}\\
L_{B}&=\mathrm{i}\sum_{i}c^{*}_{\alpha i}\partial_{t}c_{\alpha i}-\sum_{\langle ij\rangle}-t(c^{*}_{\alpha i}c_{\alpha j}+c^{*}_{\alpha j}c_{\alpha i})+\sum_{i}(-)^{i}B_{i}(t)c^{\dagger}_{i}\sigma c_{i}
\end{align*}
其中费米子拉格朗日量是二次型的。

如果我们先积掉费米子，就得到
\begin{equation*}
Z=\int \mathcal{D}\pmb{B}\;\mathrm{e}^{\mathrm{i}\int \mathrm{d}t\,L_{\mathrm{eff}}(\pmb{B})}
\end{equation*}
其中 $L_{\mathrm{eff}}(\pmb{B})$ 是有效拉格朗日量。由于 $\pmb{B}$ 与电子自旋耦合，它描述自旋波涨落。

我们先来研究 $L_{\mathrm{eff}}(\pmb{B})$ 的一些基本性质。假定 $\pmb{B}$ 在时空中为常数，则
\begin{equation*}
L_{\mathrm{eff}}(\pmb{B})=-2\sum^{\prime}_{\pmb{k}}(-E_{\pmb{k}})-\frac{3}{2U}N_{\mathrm{site}}\pmb{B}^{2}=-E_{\mathrm{mean}}(\pmb{B})
\end{equation*}
是平均场能量的负值，其中 $E_{\pmb{k}}=\sqrt{\xi^{2}_{\pmb{k}}+|\pmb{B}|^{2}}$。对 $E_{\mathrm{mean}}(\pmb{B})$ 求极小，我们得到能隙方程
\begin{equation*}
\frac{3}{2}-UN^{-1}_{\mathrm{site}}\sum^{\prime}_{\pmb{k}}\frac{1}{E_{\pmb{k}}}=0 \tag{5.6.2}
\end{equation*}
它确定 $\pmb{B}_{0}$ 的平均场值。若 $\pmb{B}_{0}\neq 0$，则平均场基态是一个 SDW 态。让我们假定情况正是如此。

显然，与不同自旋取向相联系的各基态是简并的。如果存在一个长波长的自旋取向涨落，那么局域地看，系统恰好处于众多可能基态的某一个之中。这对应于一个低能涨落。为了研究这类低能涨落，我们引入
\begin{equation*}
\pmb{B}_{\pmb{i}}(t)=|\pmb{B}_{\pmb{i}}(t)|\,\pmb{n}_{\pmb{i}}(t),\qquad |\pmb{n}_{\pmb{i}}(t)|=1
\end{equation*}
并把 $|\pmb{B}_{\pmb{i}}(t)|$ 替换为 $|\pmb{B}_{0}|$。这就给出了一个低能有效拉格朗日量 $L_{\sigma}(\pmb{n})=L_{\mathrm{eff}}(\pmb{n}\pmb{B}_{0})$。在长波极限下，有效拉格朗日量可以写成
\begin{equation*}
L_{\sigma}(\pmb{n})=\int \mathrm{d}^{d}\pmb{x}\;\frac{1}{2g}\left((\partial_{t}\pmb{n})^{2}-v^{2}(\partial_{\pmb{x}}\pmb{n})^{2}\right)+\cdots \tag{5.6.3}
\end{equation*}
由"⋯"标示的项是高阶导数项。自旋旋转对称性禁止可能出现的势能项（或质量项）$V(\pmb{n})$，并保证自旋涨落没有能隙。这样的有效拉格朗日量称为 $O(3)$ 非线性 $\sigma$ 模型。

我们注意到，有效拉格朗日量是按 $\partial_{\mu}\pmb{n}$ 的幂次展开的。然而，当 $\pmb{B}=\pmb{n}B_{0}$ 时，式 (5.6.1) 中的 $L_{B}$ 直接依赖于 $\pmb{n}$，怎样才能得到按 $\partial_{\mu}\pmb{n}$ 幂次的展开并不显然。下面我们讨论一个做到这一点的技巧。这一技巧使我们能够直接计算非线性 $\sigma$ 模型中的 $g$ 和 $v$。

我们引入 $\psi_{i}=U_{i}c_{i}$，其中
\begin{align*}
U_{i}&=\begin{pmatrix} z^{*}_{1i} & z^{*}_{2i}\\ -z_{2i} & z_{1i} \end{pmatrix}, & \psi_{i}&=\binom{\psi_{1i}}{\psi_{2i}}\\
z_{i}&=\binom{z_{1i}}{z_{2i}}, & z^{\dagger}_{i}z_{i}&=1,\qquad z^{\dagger}_{i}\sigma z_{i}=\pmb{n}_{i}
\end{align*}
这里 $\psi_{1i}$ 对应沿 $\pmb{n}_{i}$ 方向的自旋，$\psi_{2i}$ 对应沿 $-\pmb{n}_{i}$ 方向的自旋。这样 $L_{B}$ 就可以改写为
\begin{equation*}
L_{B}=\mathrm{i}\sum_{\pmb{i}}\psi^{*}_{\pmb{i}}\left(\partial_{t}+U_{\pmb{i}}\partial_{t}U^{\dagger}_{\pmb{i}}\right)\psi_{\pmb{i}}-\sum_{\langle ij\rangle}-t\left(\psi^{*}_{\pmb{i}}u_{ij}\psi_{\pmb{j}}+h.c.\right)+B_{0}\sum_{\pmb{i}}(-)^{\pmb{i}}\psi^{\dagger}_{\pmb{i}}\sigma^{z}\psi_{\pmb{i}}
\end{equation*}
其中
\begin{equation*}
u_{ij}=1+\begin{pmatrix}\; z^{\dagger}_{i}z_{j}-1 & \frac{1}{2}(z_{i}-z_{j})^{\dagger}\mathrm{i}\sigma^{y}(z^{*}_{i}+z^{*}_{j})\;\\[2pt]\; -\frac{1}{2}(z_{i}-z_{j})^{\top}\mathrm{i}\sigma^{y}(z_{i}+z_{j}) & z^{\dagger}_{j}z_{i}-1\; \end{pmatrix}
\end{equation*}
我们看到，在新基下 $B_{0}$ 项不依赖于 $\pmb{n}$；在新基下，$L_{B}$ 对 $\pmb{n}$ 的依赖是通过 $\pmb{n}$ 的导数进来的。

考虑 $\pmb{n}=\hat{\pmb{z}}$ 附近的小涨落
\begin{equation*}
\pmb{n}_{\pmb{i}}=\hat{\pmb{z}}+\mathrm{Re}\,\phi_{\pmb{i}}\,\hat{\pmb{x}}+\mathrm{Im}\,\phi_{\pmb{i}}\,\hat{\pmb{y}} \tag{5.6.4}
\end{equation*}
以及相应的旋量场
\begin{equation*}
z_{i}=\binom{1-|\phi_{i}|^{2}/8}{\phi_{i}/2}+O(\phi^{3}_{i})
\end{equation*}
再进一步假定 $\phi$ 为实数，$u_{ij}$ 便可以简化为
\begin{equation*}
u_{ij}=1+\begin{pmatrix}\; -\frac{1}{8}(\phi_{i}-\phi_{j})^{2} & \frac{1}{2}(\phi_{j}-\phi_{i})\;\\[2pt]\; -\frac{1}{2}(\phi_{j}-\phi_{i}) & -\frac{1}{8}(\phi_{i}-\phi_{j})^{2}\; \end{pmatrix}=\mathrm{e}^{\mathrm{i}\frac{\phi_{i}-\phi_{j}}{2}\sigma^{y}}
\end{equation*}
同样，$U_{i}\partial_{t}U^{\dagger}_{i}$ 可以简化为
\begin{equation*}
U_{i}\partial_{t}U^{\dagger}_{i}=\begin{pmatrix}\; 0 & -\frac{1}{2}\partial_{t}\phi_{i}\;\\[2pt]\; \frac{1}{2}\partial_{t}\phi_{i} & 0\; \end{pmatrix}
\end{equation*}
我们注意到，费米子与自旋涨落的耦合具有与 $SU(2)$ 规范场耦合的形式。对于小自旋涨落，$L_{B}$ 可以写成
\begin{equation*}
L_{B}=\mathrm{i}\sum_{i}\psi^{*}_{i}\partial_{t}\psi_{i}-H_{0}-H_{I}
\end{equation*}
其中
\begin{equation*}
H_{0}+H_{I}=-t\sum_{\langle ij\rangle}\left(\psi^{*}_{i}\,\mathrm{e}^{\mathrm{i}a_{ij}}\psi_{j}+h.c.\right)+B_{0}\sum_{i}(-)^{i}\psi^{*}_{i}\sigma^{z}\psi_{i}+\sum_{i}\psi^{*}_{i}a_{0}(\pmb{i})\psi_{i},
\end{equation*}
这里 $a_{0}(\pmb{i})=-\frac{1}{2}\partial_{t}\phi_{i}\sigma^{y}$，而 $a_{ij}=\frac{1}{2}(\phi_{i}-\phi_{j})\sigma^{y}$。

若 $a_{0}(\pmb{i})$ 与 $a_{ij}$ 在空间中为常数，则哈密顿量 $H_{0}+H_{I}$ 可以在动量空间中求解。设 $a_{0}(\pmb{i})=-\frac{1}{2}B_{y}\sigma^{y}$，$a_{ij}=\frac{1}{2}\pmb{a}\cdot(\pmb{i}-\pmb{j})\sigma^{y}$。注意 $B_{y}=\partial_{t}\phi$，$\pmb{a}=\partial_{i}\phi_{i}$。对于小的 $B_{y}$ 和 $\pmb{a}$，基态能量具有如下形式
\begin{equation*}
E_{0}(B_{y},\pmb{a})=-C_{1}B^{2}_{y}(\pmb{i})+C_{2}\pmb{a}^{2} \tag{5.6.5}
\end{equation*}
$a^{2}_{0}(\pmb{i})$ 项正是有效拉格朗日量式 (5.6.3) 中的 $(\partial_{t}\pmb{n})^{2}$ 项，即 $C_{1}a^{2}_{0}(\pmb{i})=\int \mathrm{d}^{d}\pmb{x}\,(\partial_{t}\phi)^{2}/2g$；而 $a^{2}_{ij}$ 项是 $(\partial_{\pmb{x}}\pmb{n})^{2}$ 项，即 $C_{2}\pmb{a}^{2}=\int \mathrm{d}^{d}\pmb{x}\,v^{2}(\partial_{\pmb{x}}\phi)^{2}/2g$。这使我们能够定出 $L_{\sigma}$ 中的 $g$ 和 $v$。我们发现
\begin{align*}
g^{-1}&=-a^{-d}N^{-1}_{\mathrm{site}}\frac{\partial^{2}E_{0}(B_{y})}{\partial B^{2}_{y}}\bigg|_{B_{y}=0},\\
E_{0}(B_{y})&=N^{-1}_{\mathrm{site}}\sum^{\prime}_{\pmb{k}}\left(-E^{g}_{+,\pmb{k}}-E^{g}_{-,\pmb{k}}\right), \tag{5.6.6}\\
E^{g}_{\pm,\pmb{k}}&=\sqrt{\left(\xi_{\pmb{k}}\pm\frac{B_{y}}{2}\right)^{2}+B^{2}_{0}}
\end{align*}
以及
\begin{align*}
\frac{v^{2}}{g}&=a^{-d}\frac{\partial^{2}E_{0}(\pmb{a})}{\partial\pmb{a}^{2}}\bigg|_{\pmb{a}=0},\\
E_{0}(\pmb{a})&=N^{-1}_{\mathrm{site}}\sum^{\prime}_{\pmb{k}}\left(-E^{v}_{+,\pmb{k}}-E^{v}_{-,\pmb{k}}\right), \tag{5.6.7}\\
E^{v}_{\pm,\pmb{k}}&=\left|\sqrt{\frac{1}{4}\left(\xi_{\pmb{k}+\pmb{a}/2}+\xi_{\pmb{k}-\pmb{a}/2}\right)^{2}+B^{2}_{0}}\pm\frac{1}{2}\left(\xi_{\pmb{k}+\pmb{a}/2}-\xi_{\pmb{k}-\pmb{a}/2}\right)\right|,
\end{align*}
其中 $a$ 是晶格常数，$\sum^{\prime}_{\pmb{k}}$ 表示对约化布里渊区求和（见图 5.20）。

\subsection*{问题 5.6.1}

证明式 (5.6.6) 与式 (5.6.7)。

\subsection{长程序的稳定性}

\begin{keypoints}
\item 非线性 $\sigma$ 模型使我们能够研究量子涨落与 SDW 态的稳定性。
\item 稳定性与拓扑项。
\end{keypoints}

经典地看，非线性 $\sigma$ 模型 (5.6.3) 的一个基态是 $\pmb{n}(\pmb{x})=\hat{\pmb{z}}$，它具有长程自旋序。小自旋波涨落由式 (5.6.4) 描述，其拉格朗日量为
\begin{equation*}
\frac{1}{2g}\left(|\partial_{t}\phi|^{2}-v^{2}|\partial_{\pmb{x}}\phi|^{2}\right)
\end{equation*}
具有线性色散 $\epsilon_{\pmb{k}}=v|\pmb{k}|$，正如 XY 模型中的涨落一样。量 $\left\langle|\phi(\pmb{x},t)-\phi(0)|^{2}\right\rangle$ 量度长程横向自旋涨落的强度。若 $\left\langle|\phi(\pmb{x},t)-\phi(0)|^{2}\right\rangle\gg 1$，则长程序不可能存在。这个问题与 XY 模型中的问题完全相同。涨落的强度由 $g$ 控制：$g$ 小则涨落弱。当 $g$ 太大时，长程自旋序将被自旋波涨落破坏。然而，当 $d\leqslant 1$ 时我们发现，无论 $g$ 多么小，$\left\langle|\phi(\pmb{x},t)-\phi(0)|^{2}\right\rangle$ 都会随 $(\pmb{x},t)\to\infty$ 而发散。因此，在 $d\leqslant 1$ 维，对任何 $g$ 值都不可能存在长程序。对 XY 模型而言，在某些条件下 $d=1$ 时仍可有代数长程序。现在的问题是：$O(3)$ 非线性 $\sigma$ 模型在 $d=1$ 时是否具有代数长程序。

%FIG{5.21}: {(1+1) 维 $O(3)$ 非线性 $\sigma$ 模型中的孤子/瞬子解。}

为了回答这个问题，让我们考虑虚时并研究下面的 (1+1) 维系统：
\begin{equation*}
\frac{1}{2g}\left(|\partial_{\tau}\pmb{n}|^{2}+v^{2}|\partial_{\pmb{x}}\pmb{n}|^{2}\right)
\end{equation*}
上述作用量有一个非平凡的驻定"路径"，它是一个瞬子解（见图 5.21）。该作用量不依赖于瞬子的大小，其值为 $S_{0}=4\pi v/g$。这是因为，对任何 $\pmb{n}(x,\tau)$ 和任何标度因子 $b$，都有 $S(\pmb{n}(x,\tau))=S(\pmb{n}(bx,b\tau))$。

由于瞬子的作用量有限，瞬子的（时空）密度总是有限的。瞬子之间的平均间隔是 $l\,\mathrm{e}^{S_{0}/2}$ 的量级，其中 $l$ 是短距离截断。瞬子的大小也是 $l\,\mathrm{e}^{S_{0}/2}$ 的量级（因为瞬子的作用量不依赖于其大小）。这样，长程自旋序就被瞬子破坏了。自旋--自旋关联是短程的，关联长度为
\begin{equation*}
\xi=l\,\mathrm{e}^{S_{0}/2}=l\,\mathrm{e}^{2\pi v/g} \tag{5.6.8}
\end{equation*}
有趣的是，$\xi$ 或 $1/\xi$ 对 $g$ 是非微扰的。

我们看到，$O(3)$ 非线性 $\sigma$ 模型在 1+1 维具有短程关联，且所有激发都有能隙。然而，这个结果只对了一半。$O(3)$ 非线性 $\sigma$ 模型可以包含一个改变模型低能性质的拓扑项。

为了理解拓扑项的起源，让我们考虑 Hubbard 模型 (5.5.1) 的大 $U$ 极限。当 $U\gg t$ 时，Hubbard 模型（半填充时）的低能性质由海森堡模型描述：
\begin{equation*}
H=\sum_{\langle ij\rangle}J\,\pmb{S}_{i}\cdot\pmb{S}_{j} \tag{5.6.9}
\end{equation*}
其中 $\pmb{S}_{i}=c^{\dagger}_{i}\frac{\sigma}{2}c_{i}$ 是自旋算符，$J=4t^{2}/U$。海森堡模型的有效拉格朗日量为
\begin{equation*}
L=\mathrm{i}\sum_{i}z^{\dagger}_{i}\dot{z}_{i}-\frac{J}{4}\sum_{\langle ij\rangle}\pmb{n}_{i}\cdot\pmb{n}_{j} \tag{5.6.10}
\end{equation*}
其中 $z=\binom{z_{1}}{z_{2}}$，$\pmb{n}=z^{\dagger}\sigma z$。现在让我们假定海森堡模型的基态是反铁磁（AF）态 $\pmb{n}_{i}=(-)^{i}\hat{\pmb{z}}$，并且我们想研究 AF 基态附近的小涨落。

无能隙自旋波涨落由
\begin{equation*}
\pmb{n}_{i}(t)=(-)^{i}\pmb{n}(x^{i},t)
\end{equation*}
描述。人们很想把上式代入拉格朗日量 (5.6.10) 以得到一个连续的有效理论。然而，这种简单化的做法得不到正确的结果。一个原因是，所得拉格朗日量只含一阶时间导数项而不含二阶导数项。人们可以说，在低能下，一阶时间导数项比二阶时间导数项更重要，把二阶时间导数项丢掉应该没有问题。然而事实证明，一阶时间导数项是一个全导数，对运动方程没有影响。我们需要把二阶时间导数项包括进来，运动方程才具有非平庸的动力学。这里的情形与我们在 3.3.3 节遇到的情况非常相似，在那里我们从相互作用玻色模型导出了 XY 模型。

为了得到二阶时间导数项，我们考虑下面更一般的涨落：
\begin{equation*}
\pmb{n}_{i}(t)=(-)^{i}\pmb{n}(x^{i},t)+\pmb{m}(x^{i},t) \tag{5.6.11}
\end{equation*}
其中 $\pmb{n}(x,t)$ 是 $(x,t)$ 的光滑函数，满足 $|\pmb{n}(x,t)|=1$；$\pmb{m}(x,t)$ 也是 $(x,t)$ 的光滑函数，但 $|\pmb{m}(x,t)|\ll 1$，且 $\pmb{m}(x,t)\cdot\pmb{n}(x,t)=0$。这里 $\pmb{n}$ 描述低能的 AF 自旋波涨落，而 $\pmb{m}$ 描述能量很高、将被积掉的铁磁涨落。由于 $\int \mathrm{d}t\,2\mathrm{i}z^{\dagger}_{i}\dot{z}_{i}$ 是 $\pmb{n}_{i}$ 所张的立体角，我们有
\begin{equation*}
\int \mathrm{d}t\;2\mathrm{i}z^{\dagger}_{i}\dot{z}_{i}=\int \mathrm{d}t\;2\mathrm{i}(-)^{i}z(x^{i},t)^{\dagger}\dot{z}(x^{i},t)+\int \mathrm{d}t\;\pmb{n}\cdot(\dot{\pmb{n}}\times\pmb{m})
\end{equation*}
其中 $\pmb{n}(x,t)=z^{\dagger}(x,t)\sigma z(x,t)$，即
\begin{equation*}
\mathrm{i}\sum_{i}z^{\dagger}_{i}\dot{z}_{i}=\frac{1}{2}\int \mathrm{d}x\;\frac{1}{2}\,\pmb{n}\cdot(\partial_{t}\pmb{n}\times\partial_{x}\pmb{n})+\int \mathrm{d}x\;\frac{1}{2a}\,\pmb{n}\cdot(\dot{\pmb{n}}\times\pmb{m})
\end{equation*}
其中 $a$ 是晶格间距。$\frac{J}{4}\sum_{\langle ij\rangle}\pmb{n}_{i}\cdot\pmb{n}_{j}$ 项诱导出一个 $m^{2}$ 项。这样，积掉 $\pmb{m}$ 之后，我们得到
\begin{equation*}
S=\int \mathrm{d}t\,\mathrm{d}x\;\frac{1}{2g}\left[(\partial_{t}\pmb{n})^{2}-v^{2}(\partial_{x}\pmb{n})^{2}\right]+\theta W \tag{5.6.12}
\end{equation*}
其中 $\theta=\pi$。把它与我们之前算得的有效作用量，即式 (5.6.3)，相比较，上式多出了一项 $\theta W$：
\begin{equation*}
W=\int \mathrm{d}t\,\mathrm{d}x\;\frac{1}{4\pi}\,\pmb{n}\cdot(\partial_{t}\pmb{n}\times\partial_{x}\pmb{n}) \tag{5.6.13}
\end{equation*}
这样的项是一个拓扑项，因为对任何小变化 $\delta\pmb{n}$ 都有 $W(\pmb{n}+\delta\pmb{n})-W(\pmb{n})=0$。\footnote{我们注意到
\begin{equation*}
W(\pmb{n}+\delta\pmb{n})-W(\pmb{n})=\int \mathrm{d}t\,\mathrm{d}x\left[\frac{1}{4\pi}\,\pmb{n}\cdot(\partial_{t}\delta\pmb{n}\times\partial_{x}\pmb{n})+\frac{1}{4\pi}\,\pmb{n}\cdot(\partial_{t}\pmb{n}\times\partial_{x}\delta\pmb{n})\right]
\end{equation*}
因为 $\pmb{n}\cdot\delta\pmb{n}=0$，且 $\partial_{t}\pmb{n}\times\partial_{x}\pmb{n}$ 平行于 $\pmb{n}$。另外，
\begin{align*}
\int \mathrm{d}t\,\mathrm{d}x\;\pmb{n}\cdot(\partial_{t}\delta\pmb{n}\times\partial_{x}\pmb{n})&=\int \mathrm{d}t\,\mathrm{d}x\;\partial_{t}\delta\pmb{n}\cdot(\partial_{x}\pmb{n}\times\pmb{n})\\
&=-\int \mathrm{d}t\,\mathrm{d}x\left[\delta\pmb{n}\cdot(\partial_{t}\partial_{x}\pmb{n}\times\pmb{n})+\delta\pmb{n}\cdot(\partial_{x}\pmb{n}\times\partial_{t}\pmb{n})\right]\\
&=-\int \mathrm{d}t\,\mathrm{d}x\;\delta\pmb{n}\cdot(\partial_{t}\partial_{x}\pmb{n}\times\pmb{n}),
\end{align*}
类似地，
\begin{equation*}
\int \mathrm{d}t\,\mathrm{d}x\;\pmb{n}\cdot(\partial_{t}\pmb{n}\times\partial_{x}\delta\pmb{n})=+\int \mathrm{d}t\,\mathrm{d}x\;\delta\pmb{n}\cdot(\partial_{t}\partial_{x}\pmb{n}\times\pmb{n}).
\end{equation*}} 因此，对能够连续变形到彼此的不同 $\pmb{n}$，$W(\pmb{n})$ 取相同的值。注意，$\pmb{n}$ 是从二维平面 $\mathcal{R}^{2}$ 到球面 $S^{2}$ 的一个映射，即 $\mathcal{R}^{2}\to S^{2}$。这个映射确实有不同的类，各类之间不能连续变形到彼此。这些不同的类由一个缠绕数——$\mathcal{R}^{2}$ 缠绕 $S^{2}$ 的次数——来表征。式 (5.6.13) 中的 $W(\pmb{n})$ 只取整数值，正是缠绕数的一个数学表述。

与 $g$ 和 $v$ 不同，式 (5.6.12) 中 $\theta=\pi$ 的值是精确且量子化的。它不接收微扰展开中高阶项的修正。为看清这一点，我们注意到，平移一个晶格间距在我们的有效理论 (5.6.12) 中诱导变换 $\pmb{n}\to-\pmb{n}$。在这样的变换下，$S$ 应当不变。由于在该变换下 $W\to -W$，只有 $\theta=0$ 和 $\pi$ 与平移一个晶格间距的对称性相容。这样，只要底层晶格哈密顿量具有平移一个晶格间距的对称性，$\theta$ 就只能取 0 和 $\pi$ 两个值。由于 $\theta$ 的量子化，式 (5.6.12) 不仅适用于海森堡模型 (5.6.9)，也适用于 Hubbard 模型 (5.5.1) 的 AF 态。我们还注意到，对自旋为 $S$ 的海森堡模型，其 AF 态同样由 $O(3)$ 非线性 $\sigma$ 模型描述，只是现在 $\theta=2\pi S$。（如果我们显式破坏这一对称性，比如说令 $J_{i,i+1}\neq J_{i+1,i+2}$，则 $\theta$ 的值便不再量子化。）

上面关于瞬子把代数长程序变成短程序的讨论，只在 $\theta=0\ \mathrm{mod}\ 2\pi$ 时有效。因此，整数自旋的海森堡模型具有短程关联和有限能隙，其基态是自旋单态（对偶数格点的链）。对半整数自旋的海森堡模型，$\theta=\pi$，上述瞬子讨论不再有效。事实上，这种情形下自旋--自旋关联具有代数长程序，即 $\langle\pmb{S}_{i}\cdot\pmb{S}_{j}\rangle\sim 1/|i-j|$。基态仍是自旋单态（对偶数格点的链），但现在存在无能隙激发。

缠绕数还给出瞬子作用量 $S_{0}$ 的一个下界，即 $S_{0}\geqslant 4\pi v|W|/g$。\footnote{假定 $v=1$。对于 $\delta\tau=\delta x$ 的小正方形 $\delta\tau\times\delta x$，我们有
\begin{align*}
\int_{\delta\tau\times\delta x}\mathrm{d}\tau\,\mathrm{d}x\,\left[(\partial_{\tau}\pmb{n})^{2}+(\partial_{x}\pmb{n})^{2}\right]&=(\partial_{\tau}\pmb{n})^{2}+(\partial_{x}\pmb{n})^{2}\\
&\geqslant 2\,|\pmb{n}\cdot(\partial_{\tau}\pmb{n}\times\partial_{x}\pmb{n})|=2\left|\int_{\delta\tau\times\delta x}\mathrm{d}\tau\,\mathrm{d}x\;\pmb{n}\cdot(\partial_{\tau}\pmb{n}\times\partial_{x}\pmb{n})\right|.
\end{align*}
由缠绕数 $(4\pi)^{-1}\int \mathrm{d}\tau\,\mathrm{d}x\;\pmb{n}\cdot(\partial_{\tau}\pmb{n}\times\partial_{x}\pmb{n})=\pm 1$ 可得 $S_{0}\geqslant 4\pi v|W|/g$（我们已把 $v$ 放了回来）。} 事实上，瞬子解饱和该下界，$S_{0}=4\pi v/g$。

\subsection{量子数与低能激发}

\begin{keypoints}
\item 从连续的非线性 $\sigma$ 模型计算低能激发的晶体动量。
\end{keypoints}

让我们假定 $d>1$，且系统具有真正的长程 AF 序。AF 基态之上的低能激发由 $O(3)$ 非线性 $\sigma$ 模型描述：
\begin{equation*}
S=\int \mathrm{d}t\,\mathrm{d}x\;\frac{1}{2g}\left[(\partial_{t}\pmb{n})^{2}-v^{2}(\partial_{x}\pmb{n})^{2}\right] \tag{5.6.14}
\end{equation*}
这里我们想用 $S$ 计算有限系统的低能激发。

首先，让我们考虑均匀涨落，即
\begin{equation*}
\pmb{n}(\pmb{x},t)=\pmb{n}_{0}(t) \tag{5.6.15}
\end{equation*}
我们已经假定晶格在每个方向都有偶数个格点，因此上述拟设 (5.6.15) 可以与 AF 序相容。$\pmb{n}_{0}$ 的有效作用量为
\begin{equation*}
S_{0}=\int \mathrm{d}t\;\frac{M}{2}\dot{\pmb{n}}^{2}_{0}
\end{equation*}
其中 $M=\mathcal{V}/g$，$\mathcal{V}$ 是空间的体积。这里 $S_{0}$ 描述一个质量为 $M$ 的粒子在单位球面上的运动。能级由"角动量"标记，它就是我们系统中的总自旋 $(S,S_{z})$：$E_{S,S_{z}}=\frac{S(S+1)}{2M}$，$S_{z}=-S,-S+1,\ldots,S$。由于晶格有偶数个格点，$S$ 总是整数。我们注意到，$E_{S,S_{z}}$ 按 $\mathcal{V}^{-1}$ 标度，远低于最低自旋波激发的能量——后者在 $d>1$ 时按 $\mathcal{V}^{-1/d}$ 标度。

现在让我们考虑激发态 $|S,S_{z}\rangle$ 的动量。朴素地想，人们预期所有 $|S,S_{z}\rangle$ 态都携带零动量，因为它们来自常数 $\pmb{n}_{0}$。然而，$\pmb{n}_{0}$ 只在把偶（奇）格点变到偶（奇）格点的那些平移下不变。因此，$|S,S_{z}\rangle$ 只在这些平移下不变。这把 $|S,S_{z}\rangle$ 的动量限制为 0 或 $\pmb{Q}=(\pi,\pi,\ldots,\pi)$。在平移一个晶格间距时，$\pmb{n}_{0}\to -\pmb{n}_{0}$。于是，在 $\pmb{n}_{0}\to -\pmb{n}_{0}$ 下为偶的态携带零动量，而在 $\pmb{n}_{0}\to -\pmb{n}_{0}$ 下为奇的态携带动量 $\pmb{Q}$。我们发现，偶 $S$ 态携带零动量，奇 $S$ 态携带动量 $\pmb{Q}$。

\subsection*{问题 5.6.2}

从式 (5.6.10) 出发计算式 (5.6.12) 中的作用量 $S$（即由 $J$ 和晶格间距 $a$ 确定 $g$ 与 $v$ 的值）。

% ---------------------------------------------------------------------------
% 块边界备注（装配说明，非译文）：
% 1. 起点：书页 239（p254）顶部 "This allows us to calculate" 是书页 238（p253）
%    段落（"…we find that, in real space，[实空间关系式]"）的延续。本文件第 2 行为
%    %JOIN%，第 3 行起即续文"这使我们能够计算"，与 ch5_d 文件尾装配备注的请求一致；
%    本块不设 %--E-TAIL-- 区（无上一块收尾内容需要另写）。
% 2. 终点：本块为第 5 章章末块。书页 249（p264）以问题 5.6.2 收束，其后即为第 6 章
%    章首（p265），本块自然收尾，无 TAIL、无跨页半句。
% 3. 蓝本问题备注：
%    a) 蓝本式 (5.6.1) 处残留 "tag{5}+tag{5.6.1}" 双重编号，已按 PNG 归为单 tag；
%    b) 脚注 38 的第二个公式（"⩾2|n·(∂_τn×∂_x n)|=…"）被 MinerU 错插到 5.6.3 节
%       式 (5.6.14) 之后的正文位置，本块已按 PNG 归位入脚注 38；
%    c) 式 (5.5.6) 首行蓝本作 (u_k^2−v_k^2)，经放大核对与 PNG 一致（低分辨率整页
%       判读易误作 v_k^2−u_k^2，以放大图为准）；
%    d) 蓝本 5.6 节要点行前的 "\-" 系节首要点列表，已转 keypoints 环境；
%       式 (5.6.6)/(5.6.7) 的编号在原书中挂于三行公式的中间行，本块 align* 的
%       \tag 位置照此排。
% 4. 术语：能隙方程、约化布里渊区与 ch5_d 用法一致；SDW、RPA 沿用本章前文译名，
%    不再括注；新增术语已追加至术语表"第5章（ch5_e）"一节。
